\scp{Distinguishing innovations}{07-04}
Before discussing why we exclude the peripheral languages in
the \tsc{Taliabu} (\crf{ch:Tal}) and \tsc{Seram-Tanimbar-Bomberai} (\crf{ch:STB}) subgroups,
it helps to summarise the main issues that identify
and distinguish \tsc{Sula-Buru} and \tsc{Ambon-Seram}.
\tsc{Sula-Buru} includes Mangole, Sanana,
Lisela, Buru, and Ambelau.
\tsc{Ambon-Seram} includes the languages of the western three
quarters of Seram, all of Ambon-Lease, languages in the Manipa Strait between Seram
and Buru, and the recently extinct Kayeli language of eastern Buru.
A summary is found in the classification outline
at the end of chapters \ref{ch:AS}--\ref{ch:Ban}.

We begin by comparing
and contrasting innovations in the historical phonology (\srf{07-04-01}),
and then examine developments in the lexicon (\srf{07-04-02}).
There is some overlap,
since some issues are lexically specific.
Issues in the morphosyntax are discussed in \srf{07-04-03}.